\section{Worum es geht}

\textbf{[SETZUNG]} Dieses Dokument entwickelt drei zusammenhängende Thesen:

\textbf{(1) Rotverschiebung als statischer Geometrie-Effekt.} In FFGFT entsteht
$z$ nicht durch Raumexpansion, sondern durch energetischen Verlust des Photons im
$\xi$-Feld entlang des Lichtwegs: $1+z=e^{\xi x}$.

\textbf{(2) Achromatizität.} Der Mechanismus ist rein geometrisch -- er streckt
alle Wellenlängen gleich. Frühere Darstellungen chromatischer Rotverschiebung im
Korpus sind als Korrektur im Korrektionsprogramm verzeichnet.

\textbf{(3) Der SI-Anker des Sektors.} Sowohl FFGFT als auch das Standardmodell
brauchen für den Übergang von natürlichen Einheiten zu messbaren SI-Größen in der
Kosmologie einen empirischen Anker. Dieser Anker ist in FFGFT das Verhältnis
$m_ec^2/k_BT_\text{CMB}$ -- eine gemessene Zahl, keine reine Geometrie. Das SM
hat das Analoge: $H_0$ in km/s/Mpc setzt $c$, $k_B$ und die Megaparsec-Definition
voraus, die ihrerseits historische SI-Konventionen sind.

\section{Die statische Rotverschiebung}

\textbf{[B]} Ein Photon verliert im statischen $\xi$-Feld Energie linear mit der
zurückgelegten Distanz:
\[
  \frac{dE}{dx} = -\xi\,E
  \;\Rightarrow\;
  E(x) = E_0\,e^{-\xi x}.
\]
Die Rotverschiebung:
\[
  \boxed{\;1+z = e^{\xi x},\qquad z\approx\xi\,x\;\text{ (kleine }z\text{)}\;}
\]
Kein Skalenfaktor $a(t)$, kein Urknall, kein expandierender Raum. Lichtlaufzeit
zu einem Objekt mit Rotverschiebung $z$:
\[
  \tau(z) = \frac{\ln(1+z)}{H_0^\text{T0}}.
\]
Die Funktion wächst unbeschränkt -- es gibt keine Urknall-Decke.

\textbf{[K]} \textbf{Vergleich mit dem Hubble-Gesetz.} Für kleine $z$: $z\approx\xi x$.
Hubble-Gesetz: $z=H_0 d/c$. Identifikation: $H_0=\xi c$. Die T0-Formulierung
nutzt die Hubble-Energie $E_H=E_0\cdot\xi^{41/4}$, woraus $H_0^\text{T0}=E_H/\hbar\approx66{,}2\,$km/s/Mpc
($-1{,}9\,\%$ gegenüber Planck). Der Exponent $41/4$ ist nicht aus $\xi$ allein
hergeleitet -- das ist der benannte offene Punkt .

\section{Achromatizität: Korrektur der früheren Darstellung}

\textbf{[B]} Der reife FFGFT-Mechanismus ist die \emph{fraktale Wegverlängerung}:
der Pfad durch die rekursive Zeitgeometrie wird um einen Faktor gestreckt, der
allein von der Distanz abhängt -- nicht von der Wellenlänge des Photons. Daher
gilt $1+z=e^{\xi x}$ für \emph{alle} Wellenlängen mit demselben Faktor.

\textbf{[K]} \textbf{Physikalischer Grund.} Ein \emph{geometrischer} Streckfaktor
wirkt auf Wellenlänge und Pulsdauer im selben Verhältnis -- genauso wie
Gravitationsrotverschiebung und die Standard-Kosmologie-Rotverschiebung
achromatisch sind. Ein wellenlängenabhängiger Term entstünde nur aus einem
dispersiven Medieneffekt, nicht aus reiner Geometrie.

\textbf{[K]} \textbf{Konsequenz.} Die beobachtete kosmologische Rotverschiebung
ist bis hoher Präzision achromatisch -- FFGFT trifft das korrekt. Die frühere
Darstellung chromatischer T0-Rotverschiebung in älteren Korpora-Dokumenten
in früheren Versionen des Altkorpus ist überholt; die
Korrektur ist als Korrektionseintrag verzeichnet.

\section{Der SI-Anker: wo die Zirkularität sitzt}

\textbf{[B]} Dies ist der methodisch wichtigste Abschnitt. FFGFT leitet $H_0$
aus $\xi$ und $E_0$ ab. Die Herleitung endet mit:
\[
  H_0^\text{T0} = \frac{E_H}{\hbar},
  \qquad E_H = E_0\cdot\xi^{41/4}.
\]
Scheinbar eine reine $\xi$-Ableitung. Aber $E_0=\sqrt{m_em_\mu}=7{,}398\,$MeV
ist eine \emph{gemessene} Größe (A130), und der Übergang zu km/s/Mpc braucht
$\hbar$ und die Megaparsec-Definition -- beides SI-Konventionen.

\textbf{[K]} \textbf{Der verbleibende Anker.} Nach dem Kürzen der $\hbar c$-Faktoren
(die in beiden Längenskalen $L_\xi$ und $\bar\lambda_e$ erscheinen) bleibt:
\[
  \boxed{\;\frac{m_ec^2}{k_BT_\text{CMB}} = 2{,}176\times10^9\;}
\]
Hier steckt $c^2$ (Masse $\to$ Energie) und $k_B$ (Kelvin $\to$ Energie) -- das
ist ein Verhältnis zweier \emph{gemessener} Größen: Elektronen-Ruheenergie zu
CMB-Temperatur. Die Brückenformel ist damit Geometrie-Form ($K_\text{geo}\cdot\xi^{41/4}$)
\textbf{mal} SI-Anker ($m_ec^2/k_BT_\text{CMB}$) -- keine reine Geometrie-Identität.

\textbf{[B]} \textbf{Wo der Exponent wechselt.} Dasselbe Verhältnis steckt im
Exponenten selbst: referenziert man auf die Planck-Skala statt auf $E_0$, wird
der Exponent $63/4=41/4+11/2$, wobei $11/2$ die SI-Skalen-Umrechnung
$\ell_P\leftrightarrow E_0$ kodiert. Der Faktor verschwindet nie -- er wechselt
nur den Ort.

\begin{center}
\renewcommand{\arraystretch}{1.3}
{\small
\begin{tabular}{lll}
\toprule
Referenzskala & Exponent & SI-Faktor explizit als \\
\midrule
$E_0$ (Elektron/CMB) & $41/4$ & Verhältnis $m_ec^2/k_BT_\text{CMB}$ \\
$\ell_P$ (Planck)    & $63/4 = 41/4+11/2$ & Zusatzexponent $11/2$ \\
\bottomrule
\end{tabular}
}
\renewcommand{\arraystretch}{1.0}
\end{center}

\section{Das SM hat dasselbe Problem}

\textbf{[B]} Das Standardmodell gibt $H_0$ in km/s/Mpc an. Diese Einheit setzt
voraus: $c$ (Lichtgeschwindigkeit in km/s), die Megaparsec-Definition ($1\,\text{Mpc}
=3{,}0857\times10^{22}\,\text{m}$, historisch aus Parallaxen), $k_B$ für
Temperaturmessungen und die SI-Reform 2019, die $c$, $\hbar$, $k_B$ und $e$
per Definition fixiert. Diese Größen sind nicht unabhängig voneinander -- die
Messkette ist selbstreferentiell.

\textbf{[B]} Konkret: $T_\text{CMB}=2{,}7255\,$K (Fixsen 2009) ist über $k_B$ mit Energie
verknüpft; $H_0\approx67{,}4\,$km/s/Mpc ist über $c$ und Mpc mit einer Länge
verknüpft; das Verhältnis $m_ec^2/k_BT_\text{CMB}$ verbindet Teilchenphysik mit
Kosmologie -- und dieses Verhältnis hat im SM keinen Ort in der Grundtheorie.
Das Standardmodell nennt es eine Messung; FFGFT versucht, es aus $\xi$ zu
verstehen, und kommt dem näher, aber nicht ganz hin (der Exponent $41/4$ fehlt).

\textbf{[S]} \textbf{Der ehrliche Vergleich.} Beide Rahmen -- $\Lambda$CDM und
FFGFT -- brauchen für die Kosmologie einen empirischen Anker, der außerhalb des
theoretischen Kerns liegt. Im SM ist das $H_0$ selbst (eine Messung, keine
Ableitung). In FFGFT ist es $m_ec^2/k_BT_\text{CMB}$ plus der noch nicht
hergeleitete Exponent $41/4$. FFGFT ist hier eine Stufe weiter -- aber noch nicht
vollständig parameterfrei.

\section{Kosmologische Entartung}

\textbf{[K]} FFGFT und $\Lambda$CDM sind auf den Standardtests
\emph{beobachtungs-entartet}: beide reproduzieren dieselben Flüsse, Winkelgrößen
und Zeitdehnungen. Die wichtigsten Punkte:

\textbf{Supernova-Zeitdilatation:} DES misst $b=1{,}003\pm0{,}011$ (White et al.\ 2024, Dark Energy Survey, MNRAS). Die fraktale
Wegverlängerung dehnt Wellenlänge \emph{und} Pulsdauer im selben Verhältnis
$\Rightarrow b=1$. Der Test schließt nur nicht-dilatierendes tired-light ($b=0$)
aus -- FFGFT sagt $b=1$ voraus und ist nicht betroffen.

\textbf{Hubble-Spannung:} In FFGFT ist $H_0$ keine Expansionsrate, sondern der
geometrische Energieverlustkoeffizient. Verschiedene Messmethoden sondieren
verschiedene Skalen; systematische Variationen sind geometrisch erwartet. Der
Widerspruch $67{,}4$ (Planck 2018) vs.\ $73{,}0\,$km/s/Mpc (SH0ES) ist in FFGFT kein fundamentales
Problem.

\textbf{[S]} \textbf{Ehrliche Lage.} Die Entartung schneidet in beide Richtungen:
die Daten \emph{passen} zu FFGFT (keine falsche Vorhersage), aber sie
\emph{erzwingen} es nicht. Der Fall für FFGFT ruht auf Sparsamkeit (kein $\Lambda$,
keine dunkle Energie) und auf der $\xi$-Ableitung der Konstanten -- nicht auf
einem einzelnen Diskriminator.

\section{Prüfskript}

\begin{itemize}[leftmargin=2em,itemsep=2pt,topsep=2pt]
  \item Numerische Verifikation: $1+z=e^{\xi x}$, Achromatizität (gleicher
    Streckfaktor für $\lambda_1\neq\lambda_2$), SI-Anker $m_ec^2/k_BT_\text{CMB}$,
    Exponenten-Tabelle $41/4$ vs.\ $63/4$:
    \texttt{python/A\_Serie\_Skripte/a265\_rotverschiebung.py}
\end{itemize}


\section{Zeichenerklärung}

Symbole die in der gesamten A-Serie verwendet werden, sind in A010 erklärt.
Spezifisch für dieses Dokument:

\begin{center}
\renewcommand{\arraystretch}{1.3}
{\small\begin{tabular}{lp{9cm}}
\toprule
Symbol & Bedeutung \\
\midrule
$z$ & Kosmologische Rotverschiebung; $1+z=\lambda_\text{obs}/\lambda_\text{emit}$ \\
$E_H = E_0\cdot\xi^{41/4}$ & Hubble-Energie in FFGFT \\
$H_0^\text{T0}\approx66{,}2\,$km/s/Mpc & FFGFT-Wert der Hubble-Konstante \\
$\tau(z) = \ln(1+z)/H_0^\text{T0}$ & Lichtlaufzeit in FFGFT \\
$m_ec^2/k_BT_\text{CMB}$ & Empirischer SI-Anker des kosmologischen Sektors \\
\bottomrule
\end{tabular}}
\renewcommand{\arraystretch}{1.0}
\end{center}

\section{Was offen bleibt}

\textbf{Der Exponent $41/4$ ist nicht aus $\xi$ allein hergeleitet.} Er ist an
den gemessenen $H_0$ kalibriert . Solange er fehlt, ist FFGFT im
Kosmologie-Sektor nicht parameterfrei.

\textbf{Der Anker $m_ec^2/k_BT_\text{CMB}$ braucht eine geometrische Erklärung.}
Warum gerade das Verhältnis von Elektronen-Ruheenergie zu CMB-Temperatur als
natürlicher Übergang auftritt, ist eine offene Strukturfrage.

\textbf{BAO und CMB-Peaks sind noch nicht aus FFGFT hergeleitet.} Die
Entartungsbehauptung setzt voraus, dass dieselben Winkelskalen reproduzierbar sind;
der Beweis im FFGFT-eigenen Formalismus steht aus.

\section{Ersetzt}

Dieses Dokument nimmt auf: Dok.~026 (statische Rotverschiebung als geometrische
Pfadstreckung, FEM-Simulation, Hubble-Energie $E_H=E_0\xi^{41/4}$, Herleitung
$H_0\approx66{,}2\,$km/s/Mpc), Dok.~280 (statische Rotverschiebung, Achromatizität
als Korrektur K6, $z\to$Zeit als $\Lambda$CDM-Übertragung, kosmologische Entartung,
kein Urknall), Dok.~279 (Casimir--$H_0$-Brückenformel, SI-Anker
$m_ec^2/k_BT_\text{CMB}$, Exponent $41/4$ vs.\ $63/4$, Ort des SI-Faktors).

\section{Verwendung von KI-Werkzeugen}

Dieses Dokument ist unter Verwendung KI-gestützter Werkzeuge entstanden.
Verantwortung für Inhalt und Fehler liegt beim Autor. Wortlaut der Regeln in A010.
